    \begin{figure}[t]
	\centering
	\begin{tikzpicture}[scale=0.7]
		%koren
		\node (K) at (0,0) [circle,fill=black,inner sep=1pt] {};
		%perviy uroven
		\node (A) at (-5,-2) [circle,fill=black,inner sep=1pt] {};
		\node (B) at (-2,-2) [circle,fill=black,inner sep=1pt] {};
		\node (C) at (-1,-2) [circle,fill=black,inner sep=1pt] {};
		\node (K2) at (0,-2) [circle,fill=black,inner sep=1pt] {};
		\node (D) at (1,-2) [circle,fill=black,inner sep=1pt] {};
		\node (E) at (2,-2) [circle,fill=black,inner sep=1pt] {};
		\node (F) at (5,-2) [circle,fill=black,inner sep=1pt] {};
		%zapolnaem promejutki melkimi tockami
        \draw[line width=0.8pt, dash pattern=on 0pt off 6pt, line cap=round] (-4.7,-2) -- (-2.3,-2);
        \draw[line width=0.8pt, dash pattern=on 0pt off 6pt, line cap=round] (-.7,-2) -- (.7,-2);
        \draw[line width=0.8pt, dash pattern=on 0pt off 6pt, line cap=round] (2.3,-2) -- (4.7,-2);
        %risuem linii mejdu kornem i pervim urovnem
        \draw[->, >=stealth, wred, thick](K) -- node[midway, left = 9pt] {\footnotesize $R_1$} (A);
        \draw[->, >=stealth, wred, thick] (K) -- (B);
        \draw[->, >=stealth, mblue, thick](K) -- node[midway, right = 0.1pt] {\footnotesize $R$} (C);
        \draw[->, >=stealth, mblue, thick] (K) -- (D);
        \draw[->, >=stealth, sgreen, thick] (K) -- (E);
        \draw[->, >=stealth, sgreen, thick](K) -- node[midway, right = 9pt] {\footnotesize $R_2$} (F);
        %versini dla 3qo urovna
        \node (A1) at (-7,-3.5) [circle,fill=black,inner sep=1pt] {};
        \node (B1) at (-6,-3.5) [circle,fill=black,inner sep=1pt] {};
        \node (C1) at (-5.5,-3.5) [circle,fill=black,inner sep=1pt] {};
        \node (D1) at (-4.5,-3.5) [circle,fill=black,inner sep=1pt] {};
        \node (E1) at (-4,-3.5) [circle,fill=black,inner sep=1pt] {};
        \node (F1) at (-3,-3.5) [circle,fill=black,inner sep=1pt] {};
		%risuem linii mejdu kornem i pervim urovnem
		\draw[->, >=stealth, wred, thick] (A) -- (A1);       
		\draw[->, >=stealth, wred, thick] (A) -- (B1);
		\draw[->, >=stealth, mblue, thick] (A) -- (C1);
		\draw[->, >=stealth, mblue, thick] (A) -- (D1);
		\draw[->, >=stealth, sgreen, thick] (A) -- (E1);
		\draw[->, >=stealth, sgreen, thick] (A) -- (F1);
		%zapolnaem promejutki melkimi tockami
		\draw[line width=0.8pt, dash pattern=on 0pt off 6pt, line cap=round] (-6.8,-3.5) -- (-6.1,-3.5);
		\draw[line width=0.8pt, dash pattern=on 0pt off 6pt, line cap=round] (-5.3,-3.5) -- (-4.3,-3.5);
		\draw[line width=0.8pt, dash pattern=on 0pt off 6pt, line cap=round] (-3.8,-3.5) -- (-3.1,-3.5);
		%novie linii zascet svoystv nasey logiki
		\draw[->, >=stealth, mblue, thick] (K) -- (E1);
		\draw[->, >=stealth, mblue, thick] (K) -- (F1);
		%versini dla 3qo urovna 2ya cast
		\node (A2) at (-2,-3.5) [circle,fill=black,inner sep=1pt] {};
		\node (B2) at (-1,-3.5) [circle,fill=black,inner sep=1pt] {};
		\node (C2) at (-.5,-3.5) [circle,fill=black,inner sep=1pt] {};
		\node (D2) at (.5,-3.5) [circle,fill=black,inner sep=1pt] {};
		\node (E2) at (1,-3.5) [circle,fill=black,inner sep=1pt] {};
		\node (F2) at (2,-3.5) [circle,fill=black,inner sep=1pt] {};
		%risuem linii mejdu kornem i pervim urovnem
		\draw[->, >=stealth, wred, thick] (K2) -- (A2);       
		\draw[->, >=stealth, wred, thick] (K2) -- (B2);
		\draw[->, >=stealth, mblue, thick] (K2) -- (C2);
		\draw[->, >=stealth, mblue, thick] (K2) -- (D2);
		\draw[->, >=stealth, sgreen, thick] (K2) -- (E2);
		\draw[->, >=stealth, sgreen, thick] (K2) -- (F2);
		%zapolnaem promejutki melkimi tockami
		\draw[line width=0.8pt, dash pattern=on 0pt off 6pt, line cap=round] (-1.8,-3.5) -- (-1.1,-3.5);
		\draw[line width=0.8pt, dash pattern=on 0pt off 6pt, line cap=round] (-.3,-3.5) -- (0.7,-3.5);
		\draw[line width=0.8pt, dash pattern=on 0pt off 6pt, line cap=round] (1.2,-3.5) -- (1.9,-3.5);
		%versini dla 3qo urovna 3ya cast
		\node (A3) at (3,-3.5) [circle,fill=black,inner sep=1pt] {};
		\node (B3) at (4,-3.5) [circle,fill=black,inner sep=1pt] {};
		\node (C3) at (4.5,-3.5) [circle,fill=black,inner sep=1pt] {};
		\node (D3) at (5.5,-3.5) [circle,fill=black,inner sep=1pt] {};
		\node (E3) at (6,-3.5) [circle,fill=black,inner sep=1pt] {};
		\node (F3) at (7,-3.5) [circle,fill=black,inner sep=1pt] {};
		%risuem linii mejdu kornem i pervim urovnem
		\draw[->, >=stealth, wred, thick] (F) -- (A3);       
		\draw[->, >=stealth, wred, thick] (F) -- (B3);
		\draw[->, >=stealth, mblue, thick] (F) -- (C3);
		\draw[->, >=stealth, mblue, thick] (F) -- (D3);
		\draw[->, >=stealth, sgreen, thick] (F) -- (E3);
		\draw[->, >=stealth, sgreen, thick] (F) -- (F3);
		%zapolnaem promejutki melkimi tockami
		\draw[line width=0.8pt, dash pattern=on 0pt off 6pt, line cap=round] (3.2,-3.5) -- (3.9,-3.5);
		\draw[line width=0.8pt, dash pattern=on 0pt off 6pt, line cap=round] (4.7,-3.5) -- (5.6,-3.5);
		\draw[line width=0.8pt, dash pattern=on 0pt off 6pt, line cap=round] (6.2,-3.5) -- (6.9,-3.5);
		%novie linii zascet svoystv nasey logiki
		\draw[->, >=stealth, mblue, thick] (K) -- (A3);
		\draw[->, >=stealth, mblue, thick] (K) -- (B3);
		
		\foreach \x in {-6.6,-3,0, 3, 6.5}
		{
			\foreach \y in {-3.7, -4, -4.3}
			{
				\fill (\x,\y) circle (0.75pt);
			}
		}
		
		
	\end{tikzpicture}